\usepackage[T1]{fontenc}
\usepackage{mathpazo}
\usepackage[scaled=0.92]{helvet}
\usepackage{microtype}
\usepackage{booktabs}
\usepackage{longtable}
\usepackage{array}
\usepackage{caption}
\usepackage{graphicx}
\usepackage[table]{xcolor}
\usepackage{eso-pic}

% Code looks the same in every block and inline, as in HTML: on the gray of the highlighted SQL
% blocks (shadecolor, in Quarto's Shaded box). Formula blocks (Excel, DAX, M: unhighlighted, so
% verbatim) sit in the same Shaded box; inline code (\aacode, wrapped around \texttt by
% filters/pdf-code.lua) is highlighted with lua-ul, which, unlike a box, still lets a long formula
% break at the end of a line.
\usepackage{framed}
\usepackage{luacolor}
\usepackage{lua-ul}
\providecolor{shadecolor}{RGB}{241,243,245}
\makeatletter
\@ifundefined{Shaded}{\newenvironment{Shaded}{\begin{snugshade}}{\end{snugshade}}}{}
\DeclareRobustCommand{\aacode}[1]{\highLight[shadecolor]{\hbox{}\nobreak\hskip1.5pt minus 1.5pt\nobreak
  #1\nobreak\hskip1.5pt minus 1.5pt\hbox{}}}
\AtBeginDocument{\@ifundefined{pdfstringdefDisableCommands}{}{\pdfstringdefDisableCommands{\let\aacode\@firstofone}}}
\makeatother

% Long Excel formulas in code blocks wrap after commas and asterisks, with a continuation mark,
% instead of running into the margin, and long lines of highlighted SQL (Highlighting) wrap at
% spaces. Deferred to \begin{document} because a package loaded later in Quarto's template
% redefines the verbatim environment.
\usepackage{fvextra}
\AtBeginDocument{%
  \renewenvironment{verbatim}{\VerbatimEnvironment\begin{Shaded}%
    \begin{Verbatim}[breaklines,breakafter={,*},breakaftersymbolpre={}]}{\end{Verbatim}\end{Shaded}}%
  \RecustomVerbatimEnvironment{Highlighting}{Verbatim}{commandchars=\\\{\},breaklines,breakaftersymbolpre={}}}

\definecolor{PrimaryBlue}{HTML}{1A5276}
\definecolor{Teal}{HTML}{17A589}
\definecolor{Coral}{HTML}{E74C3C}
\definecolor{Amber}{HTML}{F39C12}
\definecolor{Ink}{HTML}{2C3E50}
\definecolor{LightGray}{HTML}{F2F3F4}
\definecolor{BorderGray}{HTML}{D5D8DC}

\color{Ink}
\setkomafont{disposition}{\normalfont\sffamily\bfseries\color{PrimaryBlue}}
\setkomafont{chapter}{\Huge\sffamily\bfseries\color{PrimaryBlue}}
\setkomafont{section}{\Large\sffamily\bfseries\color{PrimaryBlue}}
\setkomafont{subsection}{\large\sffamily\bfseries\color{PrimaryBlue}}
\setkomafont{captionlabel}{\sffamily\bfseries\color{PrimaryBlue}}
\setkomafont{caption}{\small}

\captionsetup[figure]{position=bottom, justification=raggedright, singlelinecheck=false}
\captionsetup[table]{position=top, justification=raggedright, singlelinecheck=false}
\setlength{\emergencystretch}{3em}
% A higher tolerance lets a line with an unbreakable inline formula be set loosely
% instead of running into the margin.
\tolerance=3000
% A long code block that breaks across pages can leave its page a few points too full (the
% Shaded box's padding: about 9pt for a long DAX block in Chapter 14); accept that silently
% rather than report it as an overfull box.
\vfuzz=10pt
% Part entries in the table of contents set their page numbers in large bold type, and a
% three-digit number is wider than KOMA's default page-number box.
\RedeclareSectionCommand[tocpagenumberwidth=2em]{part}
\renewcommand{\arraystretch}{1.2}

\makeatletter
\renewcommand{\maketitle}{%
  \begin{titlepage}
    \thispagestyle{empty}
    \IfFileExists{visuals/cover/cover.pdf}{%
      \AddToShipoutPictureBG*{%
        \AtPageLowerLeft{%
          \includegraphics[width=\paperwidth,height=\paperheight]{visuals/cover/cover.pdf}%
        }%
      }%
      \null
    }{%
      \centering
      \vspace*{0.28\textheight}
      {\Huge\sffamily\bfseries\color{PrimaryBlue}\@title\par}
      \vspace{1.5rem}
      {\Large\sffamily\@author\par}
      \vfill
    }%
  \end{titlepage}
  \clearpage
}
\makeatother

% The title page (recto, after the cover and a blank page) and the copyright page (its verso), which
% filters/book-pages.lua writes from front-matter/_book-pages.qmd between the cover and the table of contents.
\newenvironment{aatitlepage}{%
  \cleardoubleemptypage
  \thispagestyle{empty}%
  \pdfbookmark[0]{Title Page}{book-title-page}%
  \begin{center}\sffamily\color{Ink}
  \vspace*{0.14\textheight}}{%
  \end{center}
  \clearpage}
\newcommand{\aabooktitle}[1]{{\fontsize{32}{38}\selectfont\bfseries\color{PrimaryBlue}#1\par}\vspace{0.7em}}
\newcommand{\aabooksubtitle}[1]{{\LARGE #1\par}\vspace{1.6em}{\color{Amber}\rule{1.25in}{1.5pt}}\par\vspace{1.6em}}
\newcommand{\aabookauthor}[1]{{\Large #1\par}\vfill}
\newcommand{\aabookedition}[1]{{\large #1\par}\vspace{0.3em}}
\newcommand{\aabookrevision}[1]{{\large #1\par}\vspace{2.5em}}
\newcommand{\aabookpublisher}[1]{{\large\bfseries #1\par}\vspace*{0.04\textheight}}
\newenvironment{aacopyrightpage}{%
  \thispagestyle{empty}%
  \pdfbookmark[0]{Copyright}{book-copyright-page}%
  \vspace*{\fill}
  \noindent\begin{minipage}{0.72\textwidth}\footnotesize\raggedright
  \setlength{\parindent}{0pt}\setlength{\parskip}{0.7\baselineskip}}{%
  \par\end{minipage}
  \clearpage}
